\documentclass[10pt,a4paper]{article}
\usepackage[T1]{fontenc}
\usepackage[utf8]{inputenc}
\usepackage[spanish]{babel}
\usepackage{lmodern}
\usepackage{microtype}
\usepackage[a4paper,margin=1.15cm,top=0.8cm,bottom=0.9cm]{geometry}
\usepackage{graphicx}
\usepackage{xcolor}
\usepackage{tikz}
\usepackage{booktabs}
\usepackage{parskip}
\usetikzlibrary{positioning}

\definecolor{sgreen}{HTML}{43A047}
\definecolor{sdarkgreen}{HTML}{2E7D32}
\definecolor{sblack}{HTML}{141414}
\definecolor{sdark}{HTML}{111111}
\definecolor{sgray}{HTML}{555555}
\definecolor{smidgray}{HTML}{9E9E9E}
\definecolor{slightgray}{HTML}{F2F2F2}

\graphicspath{{figures/}{brand/}}

\pagestyle{empty}
\setlength{\parindent}{0pt}
\setlength{\parskip}{0pt}

% ---------------------------------------------------------------------------
% Suyana brief macros — match the dashboard-style product-sheet layout.
% ---------------------------------------------------------------------------

% Dark header band: brand wordmark + product title + subtitle line.
\newcommand{\suyanaheader}[2]{%
  \noindent\begin{tikzpicture}[baseline=(current bounding box.north)]
    \node[fill=sdark, rounded corners=10pt, minimum width=\textwidth,
          inner xsep=18pt, inner ysep=11pt, anchor=north west] (h) at (0,0) {%
      \parbox{\dimexpr\textwidth-40pt}{%
        \color{white}%
        {\huge\textbf{Suyana}}\hfill{\Large\textbf{#1}}\\[3pt]%
        {\color{smidgray}\small\textsc{Hoja de producto}}\hfill{\color{smidgray}\small #2}%
      }%
    };
  \end{tikzpicture}\par
}

% A single stat pill — green uppercase label + white value. Used inside \suyanastats.
\newcommand{\statpill}[2]{%
  {\color{sgreen}\scriptsize\textbf{\MakeUppercase{#1}}}\\[3pt]%
  {\color{white}\normalsize #2}%
}

% Stats strip: 6 pills inside a continuation of the dark band.
\newcommand{\suyanastats}[6]{%
  \par\vspace{5pt}\noindent\begin{tikzpicture}[baseline=(current bounding box.north)]
    \node[fill=sdark, rounded corners=10pt, minimum width=\textwidth,
          inner xsep=12pt, inner ysep=8pt, anchor=north west] {%
      \parbox{\dimexpr\textwidth-28pt}{%
        \begin{minipage}[t]{0.16\linewidth}\centering #1\end{minipage}\hfill
        \begin{minipage}[t]{0.16\linewidth}\centering #2\end{minipage}\hfill
        \begin{minipage}[t]{0.16\linewidth}\centering #3\end{minipage}\hfill
        \begin{minipage}[t]{0.16\linewidth}\centering #4\end{minipage}\hfill
        \begin{minipage}[t]{0.16\linewidth}\centering #5\end{minipage}\hfill
        \begin{minipage}[t]{0.16\linewidth}\centering #6\end{minipage}%
      }%
    };
  \end{tikzpicture}\par\vspace{6pt}
}

% Big-number tile for a headline figure (AAL, prima, etc.). Use in a row.
\newcommand{\bignumtile}[2]{%
  \begin{tikzpicture}[baseline=(current bounding box.north)]
    \node[fill=white, draw=slightgray, line width=0.8pt, rounded corners=8pt,
          inner xsep=10pt, inner ysep=8pt, text width=\linewidth-22pt,
          minimum height=48pt, anchor=north west, align=center] {%
      {\fontsize{19}{22}\selectfont\bfseries\color{sgreen} #1}\\[2pt]%
      {\scriptsize\color{sgray} #2}%
    };
  \end{tikzpicture}%
}

% Content panel: green uppercase title + slightgray rounded body.
\newcommand{\suyanapanel}[2]{%
  \begin{tikzpicture}[baseline=(current bounding box.north)]
    \node[fill=slightgray, rounded corners=8pt,
          inner xsep=9pt, inner ysep=7pt, text width=\linewidth-20pt,
          anchor=north west, align=left] {%
      {\color{sgreen}\bfseries\scriptsize\MakeUppercase{#1}}\\[3pt]%
      {\color{sblack}\footnotesize #2}%
    };
  \end{tikzpicture}%
}

% Fixed-height plot block. Keeps column bottoms aligned even when captions vary.
\newcommand{\briefplot}[2]{%
  \begin{center}
    \includegraphics[width=\linewidth,height=0.115\textheight,keepaspectratio]{#1}%
  \end{center}
  \vspace{-4pt}%
  {\scriptsize\color{sgray}#2\par}%
}

% Numbered FAQ item — green circle with the number, then a content panel.
\newcommand{\faqitem}[3]{%
  \noindent\begin{minipage}{\linewidth}
    \tikz[baseline=(n.base)]\node[circle, fill=sgreen, text=white,
      font=\bfseries\scriptsize, minimum size=15pt, inner sep=0pt] (n) {#1};%
    \hspace{5pt}%
    \parbox[t]{\dimexpr\linewidth-28pt\relax}{\bfseries\scriptsize #2}%
    \par\vspace{2pt}%
    \begin{tikzpicture}[baseline=(current bounding box.north)]
      \node[fill=slightgray, rounded corners=8pt,
          inner xsep=7pt, inner ysep=5pt, text width=\linewidth-16pt,
          anchor=north west] {%
        {\color{sblack}\scriptsize #3}%
      };
    \end{tikzpicture}%
  \end{minipage}%
}

% Bottom call-to-action strip: green band, centered headline + contact.
\newcommand{\suyanacta}[2]{%
  \noindent\begin{tikzpicture}[baseline=(current bounding box.north)]
    \node[fill=sgreen, text=white, rounded corners=10pt,
          inner xsep=18pt, inner ysep=10pt,
          minimum width=\textwidth, anchor=north west] {%
      \parbox{\dimexpr\textwidth-40pt}{%
        \centering\bfseries\large #1\\[3pt]%
        \normalfont\small #2%
      }%
    };
  \end{tikzpicture}\par
}

% Footer disclaimer — small, gray, centered.
\newcommand{\suyanafooter}[1]{%
  \par\vspace{4pt}{\centering\color{smidgray}\tiny #1\par}%
}

\begin{document}
% ─────────────────────────────────────────────────────────────────────────────
% PÁGINA 1 — PRODUCTO Y MECÁNICA
% ─────────────────────────────────────────────────────────────────────────────

\suyanaheader{Producto Paramétrico de Sequía}{Argentina · Marzo 2026}

\suyanastats
  {\statpill{Región}{Sur · Norte}}
  {\statpill{Cultivo}{Trigo}}
  {\statpill{Ventana}{15 Sep – 14 Nov}}
  {\statpill{Notional}{USD 1,2M}}
  {\statpill{Climatología}{1996–2025}}
  {\statpill{Tipo}{Paramétrico}}

\vspace{5pt}

\begin{minipage}[t]{0.32\textwidth}
  \bignumtile{USD 0,10M}{Pérdida media anual histórica}
\end{minipage}\hfill
\begin{minipage}[t]{0.32\textwidth}
  \bignumtile{USD 0,52M}{Prima anual técnica}
\end{minipage}\hfill
\begin{minipage}[t]{0.32\textwidth}
  \bignumtile{USD 1,20M}{Pago máximo posible}
\end{minipage}

\par\vspace{5pt}

\begin{minipage}[t]{0.48\textwidth}

  \suyanapanel{Qué cubre este producto}{%
    Este seguro te protege cuando la sequía afecta el suelo durante la
    ventana crítica del trigo. Si el índice supera el umbral, recibís el
    pago automáticamente, sin peritaje ni espera.%
  }

  \par\vspace{6pt}

  \suyanapanel{Cuándo se activa el pago}{%
    El producto monitorea la humedad del suelo cada día de la temporada.
    Hay tres niveles de activación: sequía severa, sequía muy severa y
    sequía catastrófica. A mayor intensidad, mayor el pago que recibís.%
  }

  \par\vspace{6pt}

  \briefplot{loss_trends.png}{Pérdidas anuales históricas entre 1996 y 2025. Los años 2004 y 2009 concentraron los eventos más severos.}

\end{minipage}\hfill
\begin{minipage}[t]{0.48\textwidth}

  \suyanapanel{Estructura de pagos}{%
    \renewcommand{\arraystretch}{1.3}
    \scriptsize\setlength{\tabcolsep}{3pt}%
    \begin{tabular}{@{}lrr@{}}
      \toprule
      \textbf{Nivel de sequía} & \textbf{\% del máximo} & \textbf{Pago (USD)} \\
      \midrule
      Severa        & 20\% & 240.000 \\
      Muy severa    & 40\% & 480.000 \\
      Catastrófica  & 100\% & 1.200.000 \\
      \bottomrule
    \end{tabular}
    \par\vspace{4pt}
    {\scriptsize\color{sgray}
      Los montos corresponden al notional total del contrato.
      Cada hectárea cubierta recibe su parte proporcional.\par
    }%
  }

  \par\vspace{6pt}

  \suyanapanel{Ejemplo concreto}{%
    En 2009, el índice de humedad del suelo cayó al nivel catastrófico
    durante la ventana del trigo. El pago generado fue de USD~1.200.000,
    acreditado en forma automática al finalizar la temporada.%
  }

  \par\vspace{6pt}

  \briefplot{aep_portfolio.png}{Frecuencia histórica de pérdidas anuales del portafolio. La línea vertical marca el promedio.}

\end{minipage}

\par\vspace{6pt}
{\color{sgreen}\bfseries\scriptsize PREGUNTAS FRECUENTES}\par\vspace{3pt}

\begin{minipage}[t]{0.48\textwidth}

  \faqitem{1}{¿Qué es un seguro paramétrico?}{%
    Un seguro que paga automáticamente cuando un índice medible supera un
    umbral acordado en el contrato.%
  }

  \vspace{3pt}

  \faqitem{2}{¿Cómo medimos la sequía?}{%
    Usamos datos satelitales de humedad del suelo provistos por ERA5, calibrados
    con 30 años de historia climática local.%
  }

  \vspace{3pt}

  \faqitem{3}{¿Cómo se calcula el pago?}{%
    Si la humedad cae bajo el umbral histórico, el monto se determina
    automáticamente según el nivel de sequía alcanzado.%
  }

\end{minipage}\hfill
\begin{minipage}[t]{0.48\textwidth}

  \faqitem{4}{¿De dónde vienen los datos?}{%
    El índice usa ERA5 para humedad del suelo y controles climáticos históricos
    del pipeline de Suyana.%
  }

  \vspace{3pt}

  \faqitem{5}{¿Cuándo recibo el dinero?}{%
    El pago se procesa automáticamente al cerrar la temporada, sin necesidad
    de presentar reclamos ni documentación adicional.%
  }

  \vspace{3pt}

  \faqitem{6}{¿Qué cubre exactamente?}{%
    Cubre la sequía durante la ventana del trigo, del 15 de septiembre al
    14 de noviembre, en las regiones Sur y Norte contratadas.%
  }

\end{minipage}

\vfill

\suyanacta{¿Querés una cotización?}{Contactá a nuestro equipo en suyana.io}

\suyanafooter{Este material fue preparado por Suyana y es información
confidencial. La información presentada no constituye una oferta de
seguro ni asesoramiento. Toda cobertura paramétrica conlleva el
riesgo de que el índice gatillo no esté perfectamente correlacionado
con la pérdida subyacente.}
\end{document}
